\documentclass[11pt]{article}
\usepackage[paperwidth=16cm,paperheight=7.6cm,margin=8mm]{geometry}
\usepackage{fontspec}
\usepackage{xeCJK}
\usepackage{xcolor}
\setCJKmainfont{LXGWWenKaiGBLite-Regular.ttf}[Path=/home/liam/texlive/2026/texmf-dist/fonts/truetype/public/lxgw-fonts/]
\newfontfamily\RodFont{LXGWWenKaiGBLite-Regular.ttf}[Path=/home/liam/texlive/2026/texmf-dist/fonts/truetype/public/lxgw-fonts/]
\pagestyle{empty}
\setlength\parindent{0pt}
\newcommand\R[1]{{\RodFont #1}}
\begin{document}
{\bfseries TeX Live 的 LXGW 字体可直接排算筹（正文与算筹同一字体）}

\medskip
Unicode 核心规范 22.2 的区分例子，实测复现：\quad
\R{𝍥𝍯𝍰} = 6708，\qquad \R{𝍮𝍦𝍰} = 678

\medskip\hrule\medskip
{\bfseries 两个码区的字形朝向与区名相反}

\medskip
\begin{tabular}{@{}lll@{}}
码位 & Unicode 区名 & 实测字形 \\
\texttt{U+1D360} & COUNTING ROD \textbf{UNIT} DIGIT ONE & \R{𝍠}\quad 横（wd 10pt，ht 4.05pt） \\
\texttt{U+1D369} & COUNTING ROD \textbf{TENS} DIGIT ONE & \R{𝍩}\quad 纵（wd 10pt，ht 8.03pt） \\
\end{tabular}

\medskip
{\small 所以「个位用纵式」要取 \textbf{TENS} 区、「十位用横式」要取 \textbf{UNIT} 区——按区名直译会得到反的结果。}
\end{document}
